\subsubsection{Wake effectiveness versus distance}
Because the wake model is the dominant epistemic uncertainty, Fig.~\ref{fig:X2} maps the net fleet saving of formation transit over the wake-effectiveness factor $\kappa$ (the tabulated drag reduction is multiplied by $\kappa$; $\kappa=1$ is the planar sheet, $\kappa\approx0.7$ the free wake, $\kappa\approx0.9$ the strip model) and the one-way transit distance, with every other element of the mission simulator unchanged. The zero contour is nearly vertical in $\kappa$ at short range and bends towards short distances as $\kappa$ grows: at \SI{2.5}{km} the saving is \nkmAtTwoFive\,Wh for $\kappa=1$ and \nkmAtTwoFiveLow\,Wh for $\kappa=0.3$ (the difference is \SI{0.4}{Wh}: the transit wake benefit is almost irrelevant at that distance), while at \SI{20}{km} it ranges from small to \nkmAtTwoZerok\,Wh ($\kappa=1$). The break-even distance on the 17-strip route is \nkmBeOneSevenOne\,km at $\kappa=1$, \nkmBeOneSevenFree\,km at $\kappa=0.7$, \nkmBeOneSevenLow\,km at $\kappa=0.3$ and \nkmBeOneSevenHigh\,km at $\kappa=1.5$; on the 14-strip route (right panel) it is \nkmBeOneFourOne\,km at $\kappa=1$, \nkmBeOneFourFree\,km at $\kappa=0.7$, exceeds \SI{20}{km} (not reached within the simulated range) at $\kappa=0.3$, and is \nkmBeOneFourHigh\,km at $\kappa=1.5$. Even doubling the wake effectiveness over the best model leaves the break-even beyond \SI{4}{km}.
\begin{figure}[H]\centering\includegraphics[width=\textwidth]{X2_kappa_distance_map}
\caption{Net fleet saving of formation transit (positive = saving) over wake effectiveness $\kappa$ and one-way transit distance for (a) the 17-strip route (S3 vs S1) and (b) the 14-strip route (S4 vs D1). Dotted lines: values of $\kappa$ corresponding to the three wake models.}\label{fig:X2}\end{figure}

\subsubsection{One-at-a-time and variance-based sensitivity}
The tornado diagram (Fig.~\ref{fig:X3}a) varies one factor at a time around the baseline of the S3-versus-S1 fleet saving (\ntorBase\,Wh at \SI{2.5}{km}): wake effectiveness $\kappa\in[0.6,1.2]$ changes it by \ntorKappaLo{} and \ntorKappaHi\,Wh, the reconfiguration overhead ($0.6\times$--$1.6\times$) by \ntorReconfLo{} and \ntorReconfHi\,Wh, the station-keeping error ($0.5\times$--$3\times$) by \ntorErrLo{} and \ntorErrHi\,Wh, the wind-speed range centre (1--\SI{6}{m/s}) by \ntorWindLo{} and \ntorWindHi\,Wh, and the transit distance (1.5--\SI{10}{km}) by \ntorDistLo{} and \ntorDistHi\,Wh. Variance-based Sobol indices over the same ranges (pick--freeze estimator with $N=250$; Fig.~\ref{fig:X3}b) give first-order indices of \nsobDistS{} for distance, \nsobKappaS{} for $\kappa$, \nsobReconfS{} for reconfiguration, \nsobErrS{} for station keeping and \nsobWindS{} for wind (sum \nsobSumS); total-order indices are \nsobDistT, \nsobKappaT, \nsobReconfT, \nsobErrT{} and \nsobWindT. Transit distance dominates, followed by wake effectiveness; the interaction between them (their product is the integrated benefit) is visible as a gap between $S_i$ and $S_{Ti}$.
\begin{figure}[H]\centering\includegraphics[width=\textwidth]{X3_tornado_sobol}
\caption{(a) Tornado diagram of the S3-vs-S1 fleet saving (orange: low setting, blue: high setting); (b) Sobol first-order and total indices.}\label{fig:X3}\end{figure}
